\section{Introduction}\label{sec:intro}

Suppose you drop a ball through a Galton board with $N$ rows of pegs. At the
first peg the ball goes left or right with probability $\frac12$. At every
later peg it repeats its previous direction with probability $p$ and switches
with probability $q=1-p$. Problem~131 of Kagey's Open Problem
Collection~\cite{kagey131} asks for the law of the bin where the ball lands,
and in particular when the middle bin is exactly as likely as an endpoint. An
endpoint is reached by a single path, the one that never switches, so it has
probability $\frac12p^{N-1}$. When $p$ is close to $1$ this path wins, and when
$p$ is smaller the middle wins. Paper~A of this series~\cite{paperA} answers the
question at every $N$. The 2 are equally likely when the expected number of
switches $T=Nq$ is close to $\log N$. More precisely, the crossing value
$z_N$ of $T$ satisfies $z_N+\frac12\log z_N=\log N+\frac12\log\frac\pi8+o(1)$.

Kagey also asks what changes on other boards. In this paper the ball can move
in more than 2 directions. Our main example is Kagey's board in 2 dimensions.
A ball moves on the square lattice $\Z^2$, one step east, north, west or south
at a time. It starts in a uniform direction, and at each step it keeps its
direction, turns left or right, or reverses. Turns and reversals are rare: a
turn to each side has probability $rT/N$ and a reversal has probability $sT/N$.
After $N$ steps, which is more likely: the corner $(N,0)$, which only the path
that never turns can reach, or the origin? We show that for every even $N\ge4$
there is exactly one value $T^*_N$ of $T$ at which the 2 are equally likely,
and that
\[
 \frac{T^*_N}{\log N}\to\min\Bigl(\frac2{2r+s},\frac1s\Bigr)
\]
(Proposition~\ref{prop:decomp} and Theorem~\ref{thm:planar-first}). When
reversals dominate ($s>2r$) the origin is reached mostly by paths that stay on
one axis, and then $sT^*_N$ is the crossing $z_{N/2}$ of Kagey's board with
$N/2$ steps, up to $O(N^{-\gamma})$ (Theorem~\ref{thm:planar-rev}). When turns
dominate ($s<2r$) the origin is reached mostly by paths that use all 4
directions.

\subsection{Face selection}\label{sec:intro-faces}

The planar result is an instance of a general principle, which is the main
subject of this paper. Let the direction be a Markov chain on $d$ states that
switches from state $i$ to state $j$ with probability $q_{ij}T/N$ at each step,
and count the visits to each state. This count vector, the
\emph{composition}, plays the role of the bin. Its support is a set $F$ of
states, which we call a \emph{face} (it indexes a face of the simplex of
visit fractions). Which face carries the most likely composition, and at which
$T$ does the mode move to another face? On Kagey's board there are 2 states,
the faces of 1 state are the 2 endpoints and the face of 2 states is the set of
interior bins.

Put $T=\tau\log N$ and let $\lambda_F$ be the rate at which the chain leaves
$F$, that is, the principal Dirichlet eigenvalue on $F$ of $-Q$, where $Q$ is
the matrix of the rates $q_{ij}$. How likely the best composition with
support $F$ is depends on 2 effects.
\begin{itemize}
\item \emph{The cost of staying.} The chain must stay inside $F$ for all $N$
steps. This has probability about $e^{-T\lambda_F}=N^{-\tau\lambda_F}$. A
larger face is harder to leave, so staying on it is cheaper.
\item \emph{The cost of spreading.} The mass of $F$ is shared among the
compositions with support $F$, and there are about $N^{\lvert F\rvert-1}$ of
them, one factor $N$ for each free coordinate. Up to powers of $\log N$, the
best single composition gets a share $N^{-(\lvert F\rvert-1)}$.
\end{itemize}
So the most likely composition with support $F$ has probability
$N^{-\Phi_\tau(F)+o(1)}$, where
\[
 \Phi_\tau(F)=\tau\lambda_F+\lvert F\rvert-1,
\]
and every mode sits on a face that minimizes $\Phi_\tau$
(Theorem~\ref{thm:first}). In words, a sticky walk prefers a well-separated
community, a set of states that is hard to leave for its size. As $\tau$ grows
the cost of staying counts for more, and the mode moves to larger faces. On
Kagey's board an endpoint has $\Phi_\tau=\tau$ and the interior has
$\Phi_\tau=1$, so the middle takes over at $\tau=1$, which is the first-order
form of the answer of Paper~A. Section~\ref{sec:worked} works out 4 states on
a cycle by hand, and Theorem~\ref{thm:hull} reads the winners off the convex
hull of the points $(\lambda_F,\lvert F\rvert-1)$.

The planar walk fits this picture. The corner is the face of the single
direction east, with exit rate $2r+s$, so it has probability about
$N^{-(2r+s)\tau}$. The origin is reached either along an axis, with exit rate
$2r$, or with all 4 directions, with exit rate $0$. Comparing the 3 exponents
gives the limit above (Section~\ref{sec:planar}).

First order fixes the crossing only up to a factor $1+o(1)$. To locate it up to
$o(1)$ we need the constant in front. Say that a face $F$ satisfies (H) if the
chain can move between any 2 states of $F$ without leaving $F$, and the start
gives $F$ positive mass. For such faces we prove a local limit theorem inside
the face (Theorem~\ref{thm:local}). It shows (Theorem~\ref{thm:facemax}) that
the best composition with support $F$, with $k=\lvert F\rvert$ states, has
probability
\[
 M_F=N^{-(k-1)}\,T^{(k-1)/2}K_F\,e^{-T\lambda_F}\bigl(1+O(1/T+T^2/N)\bigr)
\]
uniformly for $1\le T\le N^{1/3}$, with an explicit constant $K_F$. The factor
$T^{(k-1)/2}$ appears because the visit fractions fluctuate on the scale
$T^{-1/2}$ in each of the $k-1$ free directions. Now let $F$ and $G$ satisfy
(H), and suppose that $F$ has $\Delta\ge1$ more states than $G$ and an exit rate
smaller by $\delta>0$. Then for large $N$ every $T\in[1,N^{1/3}]$ with $M_F=M_G$ satisfies
\[
 \delta T+\frac\Delta2\log T=\Delta\log N-\log\frac{K_F}{K_G}
 +O\Bigl(\frac1{\log N}\Bigr)
\]
(Theorem~\ref{thm:crossing}). So the coefficient $\frac12$ of $\log\log N$ that
Papers~A and~B found holds for every such pair, whatever the graph, the rates,
the sizes or the start, and without reversibility. For reversible faces $K_F$
involves the square root of a weighted count of spanning trees
(Proposition~\ref{prop:constant}), and the constants of Papers~A and~B are short
corollaries (Corollaries~\ref{cor:paperA} and~\ref{cor:paperB}). Without (H)
the coefficient can change. On a directed 3-cycle it is $1$ at one crossing and
$0$ at the next (Proposition~\ref{prop:threecycle}). Section~\ref{sec:examples}
computes all the constants on the 4-cycle, shows that the start can decide which
face wins, and lists the winners on other graphs.

The paper closes with 2 applications, and nothing before them depends on them.
Section~\ref{sec:sticky} applies face selection to sticky priors for hidden
Markov models, and Section~\ref{sec:fisher} asks how much of the Fisher
information about the switching rate is left when we only see the bin of
Kagey's board. Section~\ref{sec:model} fixes the notation,
Section~\ref{sec:verify} describes the computer checks and open questions, and
the longer proofs are in the appendices.

\subsection{The series}\label{sec:series}

This is the third of 4 papers on Problem~131, \emph{The middle and the ends
I--IV}. We call them Papers~A to~D, so that Paper~A~\cite{paperA} is part~I,
Paper~B~\cite{paperB} is part~II, this paper is part~III and
Paper~D~\cite{paperD} is part~IV. They share the following notation.
\begin{itemize}[leftmargin=1.5em,itemsep=1pt,topsep=3pt]
\item Kagey's board has $N$ steps (his row $n$ has $N=n-1$). The first step is
left or right with probability $\frac12$, and each later step repeats the
previous one with probability $p$. We write $q=1-p$ and $u=q/p$, the odds of a
switch.
\item $P_N(k)$ is the probability of $k$ right steps, that is, that the ball
lands in bin $k$. The endpoints have $P_N(0)=P_N(N)=p^{N-1}/2$.
\item $L=\log N$, and $T=Nq$ is the expected number of switches (up to the
first step).
\item The crossing is the value of $p$ (or $T$) at which a bin is exactly as
likely as an endpoint. For the middle bin it is $p_N$, and $z_N=N(1-p_N)$.
\item With $d$ states, a face $F$ is a nonempty set of states, and the vertices
are the faces of size $1$. $\lambda_F$ is the principal Dirichlet eigenvalue of
$-Q$ on $F$, and $K_F$ is the second-order constant.
\end{itemize}
Paper~A solves Kagey's board: uniform Bessel bounds bracket $z_N$ for every
$N\ge2$ and give $z_N+\frac12\log z_N=L+c_0+1/(8z_N)+R_N$ with
$c_0=\frac12\log(\pi/8)$ and an explicit $R_N$. Paper~B lets the ball switch to
each of $d-1$ other directions at the same rate; there every face ties at first
order, and constants $a_k$ break the tie at second order. Paper~D asks what
memory changes. The present paper treats every rate matrix $Q$. It explains
which face carries the mode and why, computes the constants $K_F$, recovers
$c_0$ and $a_k$ as special cases, and applies all this to the planar walk.
Figure~\ref{fig:program} compares the crossings of the 4 papers.

\begin{figure}[t]
\centering
\includegraphics[width=\textwidth]{program_C}
\caption{(a) The crossing on each model's switching scale against $L=\log N$
for Papers~A to~D. The scale is $T$ for Papers~A to~C, $x=n\varepsilon$ for the
elephant walk and the expected number of switches for the aging walk of
Paper~D. The bold curve solves the first crossing equation of
Remark~\ref{rem:planar-turns} for the planar walk with $r=s=1$, without its
$o(1)$ term. The curve A is $z_N$, and the curves B are the face crossings of
Paper~B.
(b) The exponent $e(\alpha)=1-\alpha+1/\alpha$ of Paper~D for runs with heavy
tails, which changes sign at the golden ratio $\varphi$.}
\label{fig:program}
\end{figure}

\subsection{Related work}\label{sec:related}

The rate $I_F$ of Section~\ref{sec:firstorder} is the Donsker--Varadhan
rate of the occupation measure of the chain killed on leaving $F$. Dupuis
and Liu give this rate explicitly for reversible jump
processes~\cite[Eq.~(3.2)]{dupuisliu2015}, and on a finite chain their
formula agrees with Proposition~\ref{prop:rate}(c). Guillin, Nectoux and
Wu~\cite{gnw2024} prove a large deviation principle for the occupation
measure of a process conditioned to stay in a set. For reversible processes
they show that the exit rate of the set is the least value of the
Donsker--Varadhan rate over laws on the set (their Cor.~2). In general they
prove an inequality between their rate, the Donsker--Varadhan rate and the
exit rate, and they conjecture equality (their (2.15)). For a finite chain
with $Q_F$ irreducible, the fact that $\lambda_F$ is the least value of $I_F$
is the variational formula of Donsker and Varadhan~\cite{dv1975}, which needs
no reversibility; Proposition~\ref{prop:rate}(a) gives a short proof in our
notation. This does not address the conjecture of Guillin, Nectoux and Wu,
which concerns their conditional rate. The quasi-ergodic law goes back to
Darroch and Seneta~\cite{darrochseneta1965,darrochseneta1967}, and Colonius and
Rasmussen~\cite{coloniusrasmussen2021} treat the reducible case. Dietz and
Sethuraman~\cite{dietz2007} study chains that switch from $i$ to $j$ at
time $n$ with probability $G(i,j)/n$, where every $G(i,j)>0$. Their chains
also make about $\log n$ switches, but spread log-uniformly in time, and their
occupation fractions have a limit law with full support on the simplex.
In~\cite{dietz2005} they prove large deviations for chains whose transition
matrices converge to a reducible matrix.

Exact laws of transition counts go back to Whittle~\cite{whittle1955},
Dawson and Good~\cite{dawsongood1957} and Goodman~\cite{goodman1958}.
Polettini~\cite{polettini2015} and Carugno, Vivo and Coghi~\cite{cvc2022}
derive them from the BEST theorem and the matrix-tree theorem when every
count grows linearly. We group the words by runs instead
(Lemma~\ref{lem:runs}), which suits a regime with $O(\log N)$ switches.
Bertini, Gabrielli and Landim~\cite{bgl2024} and Landim~\cite{landim2023}
expand the level-two rate functions of finite chains whose rates depend on a
parameter. They describe the time scales of such chains but do not ask where
the mode of the composition sits. The sticky HDP-HMM of Fox, Sudderth, Jordan
and Willsky~\cite{fox2011} adds a self-transition bias to the hierarchical
Dirichlet process hidden Markov model.

On the lattice the correlated random walk goes back to
Gillis~\cite{gillis1955}. Chen and Renshaw~\cite{chenrenshaw1992} treat the
Gillis--Domb--Fisher walk, and in~\cite{chenrenshaw1994} they show that the
characteristic function of any correlated random walk satisfies a
recurrence. So in principle the law of our planar walk is known. Marris and
Giuggioli~\cite{marris2024} give lattice propagators for walks with forward,
backward and lateral steps. In the continuum, Santra, Basu and
Sabhapandit~\cite{santra2020} find the marginals of a run-and-tumble particle
with 4 directions, and Angelani~\cite{angelani2022} treats switching among 4
orthogonal directions. There the corners carry atoms and the origin only a
density, so the comparison of a corner with the origin is a lattice
question. Smiga, Radaelli, Binder and Landi~\cite{smiga2023} compute the
Fisher information of the empirical distribution under a Gaussian
approximation. Their leading term vanishes for the symmetric walk of
Paper~A, whose stationary law does not depend on the switching rate. We did
not find the face-selection question or the planar crossing of the corner
and the origin in the literature.
